\documentclass[10pt,a4paper]{amsart}
\usepackage{amsmath,amssymb,amsthm,graphicx}
\usepackage{verbatim}
\newenvironment{reference}{\comment}{\endcomment}
%\newenvironment{reference}{}{}
\newenvironment{slogan}{\comment}{\endcomment}
\newenvironment{history}{\comment}{\endcomment}

% For commutative diagrams we use Xy-pic
\usepackage[all]{xy}

% We use 2cell for 2-commutative diagrams.
\xyoption{2cell}
% Theorem environments.
%
\theoremstyle{plain}
\newtheorem{theorem}[subsection]{Theorem}
\newtheorem{proposition}[subsection]{Proposition}
\newtheorem{lemma}[subsection]{Lemma}

\theoremstyle{definition}
\newtheorem{definition}[subsection]{Definition}
\newtheorem{example}[subsection]{Example}
\newtheorem{exercise}[subsection]{Exercise}
\newtheorem{situation}[subsection]{Situation}

\theoremstyle{remark}
\newtheorem{remark}[subsection]{Remark}
\newtheorem{remarks}[subsection]{Remarks}

\numberwithin{equation}{subsection}

% Macros
%
\def\lim{\mathop{\mathrm{lim}}\nolimits}
\def\colim{\mathop{\mathrm{colim}}\nolimits}
\def\Spec{\mathop{\mathrm{Spec}}}
\def\Hom{\mathop{\mathrm{Hom}}\nolimits}
\def\Ext{\mathop{\mathrm{Ext}}\nolimits}
\def\SheafHom{\mathop{\mathcal{H}\!\mathit{om}}\nolimits}
\def\SheafExt{\mathop{\mathcal{E}\!\mathit{xt}}\nolimits}
\def\Sch{\mathit{Sch}}
\def\Mor{\mathop{\mathrm{Mor}}\nolimits}
\def\Ob{\mathop{\mathrm{Ob}}\nolimits}
\def\Sh{\mathop{\mathit{Sh}}\nolimits}
\def\NL{\mathop{N\!L}\nolimits}
\def\CH{\mathop{\mathrm{CH}}\nolimits}
\def\proetale{{pro\text{-}\acute{e}tale}}
\def\etale{{\acute{e}tale}}
\def\QCoh{\mathit{QCoh}}
\def\Ker{\mathop{\mathrm{Ker}}}
\def\Im{\mathop{\mathrm{Im}}}
\def\Coker{\mathop{\mathrm{Coker}}}
\def\Coim{\mathop{\mathrm{Coim}}}

% Boxtimes
%
\DeclareMathSymbol{\boxtimes}{\mathbin}{AMSa}{"02}

%
% Macros for moduli stacks/spaces
%
\def\QCohstack{\mathcal{QC}\!\mathit{oh}}
\def\Cohstack{\mathcal{C}\!\mathit{oh}}
\def\Spacesstack{\mathcal{S}\!\mathit{paces}}
\def\Quotfunctor{\mathrm{Quot}}
\def\Hilbfunctor{\mathrm{Hilb}}
\def\Curvesstack{\mathcal{C}\!\mathit{urves}}
\def\Polarizedstack{\mathcal{P}\!\mathit{olarized}}
\def\Complexesstack{\mathcal{C}\!\mathit{omplexes}}
% \Pic is the operator that assigns to X its picard group, usage \Pic(X)
% \Picardstack_{X/B} denotes the Picard stack of X over B
% \Picardfunctor_{X/B} denotes the Picard functor of X over B
\def\Pic{\mathop{\mathrm{Pic}}\nolimits}
\def\Picardstack{\mathcal{P}\!\mathit{ic}}
\def\Picardfunctor{\mathrm{Pic}}
\def\Deformationcategory{\mathcal{D}\!\mathit{ef}}
\usepackage{cleveref}
\crefname{originalreference}{original reference}{original references}
\Crefname{originalreference}{Original reference}{Original references}
\makeatletter
\newlabel{lemma-integral-sections}{{0BH2}{source}{}{sourcepointer-1}{}}
\newlabel{lemma-integral-sections@cref}{{[originalreference][0][]0BH2}{[1][0][]source}{}{}{}}
\makeatother
\crefname{originalreference}{original reference}{original references}
\Crefname{originalreference}{Original reference}{Original references}
\makeatletter
\newlabel{lemma-normal-integral-cover-by-affines}{{0AYH}{source}{}{sourcepointer-2}{}}
\newlabel{lemma-normal-integral-cover-by-affines@cref}{{[originalreference][0][]0AYH}{[1][0][]source}{}{}{}}
\makeatother
\crefname{originallemma}{original lemma}{original lemmas}
\Crefname{originallemma}{Original lemma}{Original lemmas}
\makeatletter
\newlabel{properties-lemma-normal-integral-sections}{{0358}{source}{}{sourcepointer-3}{}}
\newlabel{properties-lemma-normal-integral-sections@cref}{{[originallemma][0][]0358}{[1][0][]source}{}{}{}}
\makeatother
\title{The ring of global functions on a normal integral algebraic space is a normal domain}
\author{The Stacks Project Authors}
\date{stacks-project:spaces-over-fields.tex:local-snapshot-20261012; GNU FDL 1.2 or later}
\begin{document}
\maketitle
\noindent\textit{Editorial note.} All delivered mathematical text is the author's own, copied from the pinned Stacks Project chapter file spaces-over-fields.tex; no statement or proof step was rewritten, reordered or removed.\par\medskip
\noindent\textit{Editorial note.} The counted proof is the author's complete argument: the reduction of normality to the integrality of elements of the fraction field, the uniqueness of local solutions over etale neighbourhoods through reducedness and the generic point condition, and the local solution by normality of the rings of sections together with the glueing.\par\medskip
\noindent\textit{Editorial note.} The label lemma-normal-integral-sections names the counted statement, and the proof environment immediately following it in the source is delivered verbatim as its proof.\par\medskip
\noindent\textit{Editorial note.} The counted statement is carried by the author's own lemma environment in the chapter source.\par\medskip
\noindent\textit{Editorial note.} Pre-open gate: python operations/agent-authoring/tools/precard\_gate.py --source spaces-over-fields.tex --label lemma-normal-integral-sections --start 538 --end 563 --local-id 28065 returned RESULT=OK (L28065 free, occupancy FREE, public interval clear). Candidate taken from sources/stacks-spaces-over-fields/candidates-halfA.json.\par\medskip
\noindent\textit{Editorial note.} The statement and proof are copied from the pinned chapter file; the cited inputs are kept as external pointers with their canonical Stacks tags.\par\medskip
\noindent\textit{Allowed standard prerequisites (editorial).}\begin{itemize}
\item The ring of global functions on an integral space is a domain. Same chapter, the lemma that the global sections of the structure sheaf on an integral algebraic space form a domain. Supplies the first half of the statement.
\item Covering by affine etale neighbourhoods with normal domains of sections. Same chapter, the lemma covering a point of a normal integral space by an affine etale neighbourhood whose ring of sections is a normal domain, and Properties, the lemma that the sections over such a neighbourhood form a normal domain. Supplies, for every point, an affine etale neighbourhood U over which the relevant fraction is integral and the sections form a normal domain.
\item Reducedness and generic points of etale neighbourhoods. Standard properties of etale morphisms of schemes: the source is reduced and its generic points map to the generic point of the target, so a section that is a nonzerodivisor globally remains a nonzerodivisor after restriction. Supplies the uniqueness of the solution of the equation b times f\_U equals a over any etale neighbourhood.
\end{itemize}\medskip
\ifdefined\hypertarget\hypertarget{sourcepointer-1}{}\fi
\noindent\textit{Source reference.} Original reference 0BH2: sources/stacks-spaces-over-fields/pinned-original-source.tex, source lines 500--500. Prerequisite pointer invoked at proof line(s) 546; it names a standard fact of the same chapter that the counted argument uses, and every step of the counted proof built on it is retained verbatim in the delivered body.\par\medskip
\ifdefined\hypertarget\hypertarget{sourcepointer-2}{}\fi
\noindent\textit{Source reference.} Original reference 0AYH: sources/stacks-spaces-over-fields/pinned-original-source.tex, source lines 517--517. Prerequisite pointer invoked at proof line(s) 555; it names a standard fact of the same chapter that the counted argument uses, and every step of the counted proof built on it is retained verbatim in the delivered body.\par\medskip
\ifdefined\hypertarget\hypertarget{sourcepointer-3}{}\fi
\noindent\textit{Source reference.} Original lemma 0358: sources/stacks-chow/tags.txt, source lines 4119--4119. Prerequisite pointer invoked at proof line(s) 559; it names a standard fact of a sibling chapter that the counted argument uses, and every step of the counted proof built on it is retained verbatim in the delivered body.\par\medskip
\begin{lemma}
\label{lemma-normal-integral-sections}
Let $S$ be a scheme. Let $X$ be a normal integral algebraic space over $S$.
Then $\Gamma(X, \mathcal{O}_X)$ is a normal domain.
\end{lemma}

\begin{proof}
Set $R = \Gamma(X, \mathcal{O}_X)$. Then $R$ is a domain by
Lemma \ref{lemma-integral-sections}.
Let $f = a/b$ be an element of the fraction field of $R$
which is integral over $R$.
For any $U \to X$ \'etale with $U$ a scheme there is at most one
$f_U \in \Gamma(U, \mathcal{O}_U)$ with $b|_U f_U = a|_U$.
Namely, $U$ is reduced and the generic points of $U$ map to
the generic point of $X$ which implies that $b|_U$ is a
nonzerodivisor.
For every $x \in |X|$ we choose $U \to X$ as in
Lemma \ref{lemma-normal-integral-cover-by-affines}.
Then there is a unique $f_U \in \Gamma(U, \mathcal{O}_U)$
with $b|_U f_U = a|_U$ because
$\Gamma(U, \mathcal{O}_U)$ is a normal domain by
Properties, Lemma \ref{properties-lemma-normal-integral-sections}.
By the uniqueness mentioned above these $f_U$
glue and define a global section $f$ of the structure
sheaf, i.e., of $R$.
\end{proof}

\end{document}
